\chapter{Inventario de archivos entregados}

\small
\begin{longtable}{p{5.3cm}p{8.0cm}}
\toprule
\textbf{Ruta} & \textbf{Propósito} \\
\midrule
\endfirsthead
\toprule
\textbf{Ruta} & \textbf{Propósito} \\
\midrule
\endhead
\path{docs/fase-1/main.tex} & Documento maestro LaTeX. \\
\path{docs/fase-1/preamble.tex} & Tipografía, colores, geometría y componentes visuales. \\
\path{docs/fase-1/sections/} & Capítulos modulares del informe. \\
\path{docs/fase-1/assets/diagrams/} & Diagramas PDF vectoriales y vistas PNG. \\
\path{docs/fase-1/assets/generated/} & Ilustración original de portada. \\
\path{docs/fase-1/assets/screenshots/} & Evidencias locales y de preparación incorporadas al informe. \\
\path{infra/terraform/} & VPC, subredes, rutas, seguridad, cómputo y salidas. \\
\path{scripts/generate-figures.py} & Generación reproducible de figuras con Matplotlib. \\
\path{scripts/validate-network.py} & Pruebas matemáticas del direccionamiento. \\
\path{scripts/test-connectivity.ps1} & Ping privado con salto SSH y registro de evidencia. \\
\path{scripts/build-report.ps1} & Generación de figuras, validación y compilación del PDF. \\
\path{evidence/} & Resultados verificables, salidas de Terraform y capturas genuinas de AWS. \\
\path{output/pdf/} & PDF final estable. \\
\href{https://github.com/DanielBarillasM/Proyecto-final_Grupo-1_Redes_Sec-10.git}{Repositorio GitHub del proyecto} & Código fuente, informe, scripts, evidencia permitida e historial de versiones. \\
\bottomrule
\end{longtable}
\normalsize

Las figuras se generaron con la interfaz orientada a objetos de Matplotlib y se exportaron como PDF vectorial y PNG de revisión. Esta metodología adopta prácticas de composición, accesibilidad y exportación documentadas en Scientific Agent Skills \cite{kassis2026skills}; el script queda incluido para que colores, etiquetas y datos puedan auditarse.

\chapter{Extractos de configuración}

\section{Definición de segmentos}

\begin{Verbatim}[fontsize=\footnotesize,frame=single]
subnets = {
  visitas        = { cidr = "10.10.0.0/26",   public = true  }
  ventas         = { cidr = "10.10.0.64/26",  public = false }
  ti             = { cidr = "10.10.0.128/27", public = false }
  datacenter     = { cidr = "10.10.0.160/28", public = false }
  administracion = { cidr = "10.10.0.176/28", public = true  }
}
\end{Verbatim}

\section{Regla ICMP de Ventas}

\begin{Verbatim}[fontsize=\footnotesize,frame=single]
ingress {
  description = "ICMP entre miembros del mismo grupo"
  from_port   = -1
  to_port     = -1
  protocol    = "icmp"
  self        = true
}
\end{Verbatim}

\section{Salidas relevantes}

\begin{Verbatim}[fontsize=\footnotesize,frame=single]
output "bastion_public_ip" {
  value = try(aws_instance.bastion[0].public_ip, null)
}

output "ventas_private_ips" {
  value = aws_instance.ventas[*].private_ip
}
\end{Verbatim}

\chapter{Procedimiento operativo abreviado}

\begin{enumerate}
  \item Autenticarse en AWS por un mecanismo local autorizado.
  \item Consultar identidad y región sin registrar secretos.
  \item Generar llave SSH temporal.
  \item Copiar y completar \code{terraform.tfvars.example}.
  \item Ejecutar \code{terraform init}, \code{fmt}, \code{validate} y \code{plan}.
  \item Revisar el plan con otro integrante.
  \item Aplicar y esperar estado saludable de las instancias.
  \item Ejecutar \code{scripts/test-connectivity.ps1} con las salidas obtenidas.
  \item Capturar VPC, subredes, rutas, reglas, instancias y resultados.
  \item Ejecutar \code{terraform destroy}.
  \item Auditar la ausencia de recursos etiquetados.
  \item Recompilar, inspeccionar el PDF y publicar el repositorio.
\end{enumerate}
